% year: תשפ"ד
% semester: ב
% moed: ג
% number: 3
% categories: func-analysis
% source: exams/מבחנים/תשפד/מועד ג תשפד סמסטר ב.pdf

\begin{question}
חקרו את הפונקציה הבאה (נק' קיצון מקומיות, תחומי עליה/ירידה, אסימפטוטות):
\[ f(x)=e^{-(x^3-3x)} \]
שרטטו את הגרף שלה.
\end{question}

\begin{hint}
$f=e^{g}$ עם $g(x)=3x-x^3$; מכיוון ש-$e^t$ עולה, $f$ עולה/יורדת בדיוק היכן ש-$g$ עולה/יורדת.
\end{hint}

\begin{solution}
\textbf{תחום הגדרה:} כל $\R$; $f$ רציפה וגזירה בכל מקום, ו-$f(x)>0$ לכל $x$ (אין חיתוך עם ציר $x$). חיתוך עם ציר $y$: $f(0)=e^0=1$.

\textbf{נגזרת:} לפי כלל השרשרת
\[ f'(x)=e^{-(x^3-3x)}\cdot(-(3x^2-3))=3(1-x^2)\,e^{-(x^3-3x)}=3(1-x)(1+x)\,e^{3x-x^3}. \]
האקספוננט חיובי, לכן סימן $f'$ הוא סימן $1-x^2$:
\begin{itemize}
  \item $(-\infty,-1)$: $f'<0$, $f$ \textbf{יורדת}.
  \item $(-1,1)$: $f'>0$, $f$ \textbf{עולה}.
  \item $(1,\infty)$: $f'<0$, $f$ \textbf{יורדת}.
\end{itemize}

\textbf{נקודות קיצון:} ב-$x=-1$ הנגזרת עוברת מ-$-$ ל-$+$: \textbf{מינימום מקומי} $f(-1)=e^{-(-1+3)}=e^{-2}\approx0.135$. ב-$x=1$ הנגזרת עוברת מ-$+$ ל-$-$: \textbf{מקסימום מקומי} $f(1)=e^{-(1-3)}=e^{2}\approx7.39$.

\textbf{אסימפטוטות:}
\begin{itemize}
  \item אנכיות: אין, $f$ רציפה על כל הישר.
  \item $x\to+\infty$: $3x-x^3=-x^3\left(1-\frac{3}{x^2}\right)\to-\infty$, לכן $f(x)\to0$: אסימפטוטה אופקית $y=0$ (הגרף מעליה).
  \item $x\to-\infty$: $3x-x^3\to+\infty$, לכן $f(x)\to+\infty$; גם $\frac{f(x)}{x}\to-\infty$ (מונה $\to+\infty$, מכנה $\to-\infty$), ולכן אין אסימפטוטה אופקית או משופעת ב-$-\infty$.
\end{itemize}

\textbf{הגרף:} מגיע מ-$+\infty$ בצד שמאל, יורד עד המינימום המקומי $(-1,e^{-2})$, עולה דרך $(0,1)$ עד המקסימום המקומי $(1,e^2)$, ואז יורד ושואף ל-$0$ (מעל ציר $x$) כאשר $x\to\infty$. הגרף כולו מעל ציר $x$.
\end{solution}
